\section{Conclusiones}
\label{sec:conclusiones}

Este trabajo construy\'o un \'indice de prioridad de revisi\'on para presas mexicanas, enteramente a partir de datos abiertos de terceros --peligro s\'ismico (GEM Foundation), lluvia extrema hist\'orica (NASA POWER) y poblaci\'on cercana (GeoNames)-- motivado por el reconocimiento p\'ublico de CONAGUA de que carece de presupuesto para inspeccionar cerca de 1{,}000 de sus presas registradas. Aplicado a las 210 presas que CONAGUA s\'i monitorea activamente, el \'indice produjo un ranking geogr\'aficamente coherente con la sismolog\'ia e hidrolog\'ia conocidas del pa\'is, combinando de forma no trivial presas de alto peligro f\'isico con presas de alta exposici\'on poblacional --tres factores que result\'aron pr\'acticamente independientes entre s\'i (Secci\'on \ref{sec:resultados})--, demostrando que la metodolog\'ia es t\'ecnicamente viable y reproducible con recursos de un solo investigador y sin financiamiento institucional.

El trabajo es expl\'icito, sin embargo, en que mide \emph{peligro y exposici\'on}, no \emph{vulnerabilidad estructural} --y en que se aplic\'o sobre el subconjunto de presas mejor vigilado del pa\'is, no sobre las que motivan el proyecto--. Cerrar esa brecha depende de una solicitud de acceso a informaci\'on p\'ublica presentada ante la Plataforma Nacional de Transparencia, pidiendo el inventario nacional completo con atributos estructurales (altura de cortina, capacidad del vertedor, a\~no de construcci\'on, estado f\'isico y fecha de \'ultima inspecci\'on). El pipeline de este trabajo --disponible como c\'odigo abierto en \url{https://github.com/mxnueel/riesgo-presas-mexico}-- est\'a preparado para incorporar esos datos en cuanto se reciban, extendiendo el an\'alisis del subconjunto de 210 presas monitoreadas al inventario completo de \(\sim\)5{,}166 presas registradas y, en la medida en que se documenten, a los \(\sim\)8{,}000 bordos no registrados.

Como trabajo futuro inmediato, adem\'as de incorporar la vulnerabilidad estructural real, quedan pendientes: reemplazar el proxy de poblaci\'on en l\'inea recta por una zona de inundaci\'on real derivada de un modelo de elevaci\'on digital; e incorporar la clasificaci\'on de riesgo ya publicada por CONAGUA para las 95 presas de alto riesgo conocidas, como una verificaci\'on directa de si el \'indice las se\~nala consistentemente como prioritarias.
